\newpage
\section{Future Work \& Summary}
\subsection{Future Work}

To address the technical constraints and empirical findings identified during system evaluation, future research and development for the H.E.R.A. platform will focus on four core objectives:

\begin{itemize}
    \item \textbf{Dynamic Device Management Integration:} Expand the system architecture across the edge firmware, MongoDB database schema, and web dashboard to support dynamic device provisioning. This feature will allow users to add, delete, and modify the operational states of new hardware devices directly through the user interface, enabling seamless scalability without requiring manual code reconfiguration or firmware recompilation.
    \item \textbf{Comprehensive AI Voice Interaction Pipeline:} Advance the communication capabilities of the H.E.R.A. assistant by transitioning from the current text-based command execution to an end-to-end speech-driven framework. Future iterations will integrate robust Speech-to-Text (STT) and Text-to-Speech (TTS) models, enabling the conversational pipeline to process direct audio input and deliver natural language voice responses to residents.
    \item \textbf{Context-Aware AI Personalization:} Enhance the machine learning infrastructure to analyze individual behavioral trends and multi-turn interaction histories for each household member. By learning routine patterns and schedules over time, the system will autonomously suggest custom automation presets and recognize specialized, user-specific voice commands tailored to individual resident preferences.
    \item \textbf{UI Synchronization and Latency Optimization:} Rectify the synchronization bottleneck observed between the physical hardware actuators and the web frontend. While Remote Procedure Call (RPC) execution at the edge level remains highly efficient (under \SI{10}{\ms}), future work will implement backend payload downsampling, introduce optimistic updates on the client side, and optimize the Server-Sent Events (SSE) data stream to eliminate the current delay and ensure immediate visual state updates.
\end{itemize}
\subsection{Summary}
The \textbf{H.E.R.A} project is an AIoT virtual assistant system that optimizes the integration between edge device control and safe artificial intelligence. Regarding hardware, the system utilizes the ESP32-S3 microcontroller (OhStem Yolo UNO) running firmware based on a multitasking FreeRTOS architecture. By decoupling IoT tasks to run concurrently, the hardware achieves ultra-low response latency ($<10$ ms) and maintains stability through automatic network connection recovery. The core highlight lies in the AI layer with a Multi-Agent architecture built on LangGraph, establishing a strict boundary between linguistic reasoning and physical execution. The LLM (Ollama/OpenRouter) is not allowed to directly manipulate hardware but only generates tool proposals ; subsequently, an independent Policy Engine strictly validates these proposals to completely eliminate the risk of hallucination before execution. Edge telemetry data is then stored at high speed using a MongoDB database (optimized for time-series data)  and instantly synchronized (via SSE) to the React Web interface ---which acts as a Digital Twin for intuitive user monitoring and control.

Detailed source code of the project: \href{https://github.com/TienHwng/MP-AI-252.git}{GitHub Link} 